\scanpage{74}
l'assertion précédente (que la section $\psi$ de $E_{U_\eta}$ sur $E_K$ provient d'une section de $E_{U_S}$ sur $E_S$, sss la section rationnelle correspondante est morphique)

Mais il s'agit ici d'un énoncé en fait purement géométrique, que j'ai envie de reformuler sous forme plus générale :

\uline{Théorème.} \deleted{\unsure{2|3}}\unsure{(?)} Soit $\unsure{S|T}$ un trait (spectre d'un anneau-jauge), $U$ un schéma relatif ``élémentaire'' sur $T$, anabélien*, $K$ le corps des fonctions de $T$, On choisit un rev. universel $\widetilde U$ de $U$, d'où une clôture alg. $\overline K$ de $K$, et on considère l'extension
\[
E_U=\pi_1(U;\widetilde U)\quad\text{de}\quad E_K=\pi_1(K,\overline K)\simeq\operatorname{Gal}(\overline K/K)
\]
par $\pi=\pi_1(U_{\overline K},\widetilde U)$. \deleted{Considérons une} \inserted{On a donc} \deleted{section pt $K$-rationnel de $U_K$,} \inserted{un carré cartésien de groupes profinis}

\uline{d'où}
\[
\diagram{\begin{tikzcd}[ampersand replacement=\&]
E_U \arrow[d,"\mathrm{surj}"'] \& E_{U_\eta} \arrow[l,"\mathrm{surj}"'] \arrow[d,"\mathrm{surj}"] \\
E_S \& E_K=\operatorname{Gal}(\overline K/K) \arrow[l,"\mathrm{surj}"]
\end{tikzcd}}
\]
où $E_S$ s'identifie au quotient de $E_K$ par le sous-groupe nor. fermé engendré par un groupe
\margin{* anabélien est inutile -- il suffit que les fibres de dim 1 des fibrations élémentaires de $U$ ne soient pas des types $(0,0)$ ou $(0,1)$ -- i.e. à $\pi_1$ nul}
\margin{73}
